%=============================================================
%  Chapter 5 --- Iteration 1 : E-Commerce Core
%=============================================================
\chapter{Implementation --- Iteration 1: E-Commerce Core}

\section*{Introduction}
The first iteration delivers a functional e-commerce foundation covering user
authentication, catalogue browsing, cart management, and order placement.
By the end of this sprint, a customer can register, browse products,
add items to the cart, and complete an order with cash-on-delivery payment.

%-------------------------------------------------------------
\section{Sprint Planning}
%-------------------------------------------------------------

\subsection{Sprint Goal}

\begin{tcolorbox}[colback=blue!5, colframe=blue!40, title=Sprint Goal --- Iteration 1]
Allow a visitor to become an authenticated customer and place their first
order end-to-end, without any manual intervention.
\end{tcolorbox}

\subsection{Sprint Backlog}

\begin{table}[H]
\centering
\caption{User stories --- Iteration 1}
\begin{tabularx}{\textwidth}{|c|X|c|c|}
\hline
\rowcolor{gray!20}
\textbf{ID} & \textbf{User Story} & \textbf{Priority} & \textbf{Points} \\
\hline
US-01 & As a visitor, I want to register with my email and receive an OTP confirmation code. & Must & 5 \\
\hline
US-02 & As a visitor, I want to log in with email and password. & Must & 3 \\
\hline
US-03 & As a customer, I want to enable two-factor authentication (TOTP). & Should & 5 \\
\hline
US-04 & As a customer, I want to reset my password via email. & Should & 3 \\
\hline
US-05 & As a visitor, I want to browse the catalogue with filters (category, price, skin type). & Must & 8 \\
\hline
US-06 & As a visitor, I want to view a detailed product page. & Must & 3 \\
\hline
US-07 & As a customer, I want to add, update, and remove items in my cart. & Must & 5 \\
\hline
US-08 & As a customer, I want to complete checkout with a shipping address. & Must & 8 \\
\hline
US-09 & As a customer, I want to view my order history. & Must & 3 \\
\hline
US-10 & As an admin, I want to create, edit, and delete products. & Must & 8 \\
\hline
\end{tabularx}
\end{table}

%-------------------------------------------------------------
\section{Design}
%-------------------------------------------------------------

\subsection{Use Case Scope}

The use cases covered by this iteration are detailed in the \textit{uc\_auth}
and \textit{uc\_shopping} diagrams presented in Chapter 2.
The main actors are: the \textbf{Visitor} (catalogue access),
the \textbf{Customer} (cart, order, profile), and the \textbf{Admin}
(product management).

\subsection{Data Model}

The entities involved in this iteration are: \texttt{User}, \texttt{Product},
\texttt{Category}, \texttt{ProductVariant}, \texttt{Cart}, \texttt{CartItem},
\texttt{Order}, and \texttt{OrderItem}, whose full class diagram is presented
in figure~\ref{fig:class_domain} of Chapter 3.

%-------------------------------------------------------------
\section{Implementation}
%-------------------------------------------------------------

\subsection{Authentication Module}

\subsubsection{Registration and OTP Verification}

Registration uses the Better Auth \texttt{emailOTP} plugin.
On form submission, the server creates the user with
\texttt{emailVerified = false} and sends a six-digit code by email.
The user is redirected to a verification page; the session is only granted
after the code is validated.

\begin{lstlisting}[language=Java, caption=Registration flow --- server side (AuthController)]
@Post('/register')
async register(@Body() dto: RegisterDto, @Res() res: Response) {
  const result = await this.authService.createUser(dto);
  if (!result.success)
    return res.status(400).json({ message: result.error });
  // Better Auth emailOTP plugin sends the code automatically
  return res.status(201).json({ message: 'Check your email.' });
}
\end{lstlisting}

\begin{figure}[H]
  \centering
  \includegraphics[width=0.85\textwidth]{figures/screen_register.png}
  \caption{Registration page --- form and OTP verification}
  \label{fig:screen_register}
\end{figure}

\subsubsection{Login and Brute-Force Protection}

The \texttt{POST /api/auth/sign-in/email} endpoint is protected by a Redis
sliding window: five failed attempts per IP/email within ten minutes trigger
a temporary lockout with a \texttt{429 Too Many Requests} response.

\begin{lstlisting}[language=Java, caption=Brute-force protection middleware]
const key = `login_attempts:${ip}:${email}`;
const attempts = await redisClient.incr(key);
if (attempts === 1) await redisClient.expire(key, 600); // 10 min window
if (attempts > 5)
  return res.status(429).json({ message: 'Too many attempts. Try later.' });
\end{lstlisting}

\begin{figure}[H]
  \centering
  \includegraphics[width=0.85\textwidth]{figures/screen_login.png}
  \caption{Login page}
  \label{fig:screen_login}
\end{figure}

\subsubsection{Two-Factor Authentication (TOTP)}

The Better Auth \texttt{twoFactor} plugin generates a TOTP secret linked to
the account and returns a QR code to scan with Google Authenticator or Authy.
Once enabled, every login requires the six-digit code in addition to the password.

\begin{figure}[H]
  \centering
  \includegraphics[width=0.75\textwidth]{figures/screen_2fa.png}
  \caption{Two-factor authentication setup --- TOTP QR Code}
  \label{fig:screen_2fa}
\end{figure}

%-------------------------------------------------------------
\subsection{Product Catalogue}
%-------------------------------------------------------------

\subsubsection{Database-Side Filtering}

Product filtering is fully delegated to PostgreSQL via Drizzle ORM,
avoiding any in-memory filtering.
Available filters are: category, price range, skin type (JSONB column),
and case-insensitive text search.

\begin{lstlisting}[language=Java, caption=Filtered query --- ProductRepository.findFiltered()]
const conditions = [eq(products.isActive, true)];
if (input.categoryId) conditions.push(eq(products.categoryId, input.categoryId));
if (input.search)     conditions.push(ilike(products.name, `%${input.search}%`));
if (input.minPrice)   conditions.push(gte(products.price, input.minPrice));
if (input.maxPrice)   conditions.push(lte(products.price, input.maxPrice));
if (input.skinType)
  conditions.push(sql`${products.skinType} @> ${JSON.stringify([input.skinType])}::jsonb`);

return db.select().from(products).where(and(...conditions));
\end{lstlisting}

\begin{figure}[H]
  \centering
  \includegraphics[width=\textwidth]{figures/screen_catalogue.png}
  \caption{Catalogue page --- product grid with sidebar filters}
  \label{fig:screen_catalogue}
\end{figure}

\subsubsection{Product Page}

The product page displays images, description, ingredients, customer reviews,
and a variant selector (size, scent).
It is server-rendered (SSR) for SEO indexing.

\begin{figure}[H]
  \centering
  \includegraphics[width=\textwidth]{figures/screen_product.png}
  \caption{Product page --- details, variants, and add-to-cart button}
  \label{fig:screen_product}
\end{figure}

%-------------------------------------------------------------
\subsection{Cart}
%-------------------------------------------------------------

The cart is persisted in the database (\texttt{cart\_items} table),
enabling session continuity across devices.
The add, update quantity, and remove operations are exposed via REST API
and consumed from the Zustand store on the client side.

\begin{figure}[H]
  \centering
  \includegraphics[width=0.85\textwidth]{figures/screen_cart.png}
  \caption{Cart page --- summary and checkout button}
  \label{fig:screen_cart}
\end{figure}

%-------------------------------------------------------------
\subsection{Checkout}
%-------------------------------------------------------------

The checkout flow is split into two steps (Shipping -> Payment) with
Zod validation on the client side and stock verification on the server side.
The order is created inside an atomic transaction: stock decrement,
order creation, and cart clearing are either all committed or all rolled back.

\begin{lstlisting}[language=Java, caption=Atomic transaction in CheckoutUseCase]
const order = await this.transactionManager.execute(async () => {
  for (const item of cart.items) {
    // validate product + decrement stock
    product.stock -= item.quantity;
    await this.productRepo.update(product);
    orderItems.push(new OrderItem(...));
  }
  await this.orderRepo.create(order);
  await this.cartRepo.clearCart(userId);
  return order;
});
\end{lstlisting}

\begin{figure}[H]
  \centering
  \includegraphics[width=\textwidth]{figures/screen_checkout.png}
  \caption{Checkout flow --- Shipping and Payment steps}
  \label{fig:screen_checkout}
\end{figure}

%-------------------------------------------------------------
\subsection{Order History}
%-------------------------------------------------------------

The \texttt{/orders} page lists all orders for the logged-in customer,
with real-time status (\textsc{pending}, \textsc{confirmed}, \textsc{shipped},
\textsc{delivered}, \textsc{cancelled}) and a line-item summary.

\begin{figure}[H]
  \centering
  \includegraphics[width=\textwidth]{figures/screen_orders.png}
  \caption{Order history page}
  \label{fig:screen_orders}
\end{figure}

%-------------------------------------------------------------
\subsection{Admin --- Product Management}
%-------------------------------------------------------------

The admin accesses a dedicated dashboard (\texttt{/admin}) protected by a
role middleware (\texttt{role: 'admin' | 'super\_admin'}).
Product management supports creation with image upload, price/stock/variant
editing, and soft deletion (\texttt{isActive = false}) without physical removal.

\begin{figure}[H]
  \centering
  \includegraphics[width=\textwidth]{figures/screen_admin_products.png}
  \caption{Admin interface --- product list and management}
  \label{fig:screen_admin_products}
\end{figure}

\begin{figure}[H]
  \centering
  \includegraphics[width=0.85\textwidth]{figures/screen_admin_create_product.png}
  \caption{Admin interface --- product creation form}
  \label{fig:screen_admin_create_product}
\end{figure}

%-------------------------------------------------------------
\section{Sprint Review}
%-------------------------------------------------------------

\subsection{Delivered Features}

\begin{table}[H]
\centering
\caption{Iteration 1 results}
\begin{tabular}{|c|l|c|}
\hline
\rowcolor{gray!20}
\textbf{ID} & \textbf{User Story} & \textbf{Status} \\
\hline
US-01 & Registration + OTP verification & \textcolor{green!60!black}{\checkmark\ Done} \\
\hline
US-02 & Email + password login & \textcolor{green!60!black}{\checkmark\ Done} \\
\hline
US-03 & TOTP two-factor authentication & \textcolor{green!60!black}{\checkmark\ Done} \\
\hline
US-04 & Password reset & \textcolor{green!60!black}{\checkmark\ Done} \\
\hline
US-05 & Catalogue with filters & \textcolor{green!60!black}{\checkmark\ Done} \\
\hline
US-06 & Detailed product page & \textcolor{green!60!black}{\checkmark\ Done} \\
\hline
US-07 & Cart management & \textcolor{green!60!black}{\checkmark\ Done} \\
\hline
US-08 & Checkout (cash on delivery) & \textcolor{green!60!black}{\checkmark\ Done} \\
\hline
US-09 & Order history & \textcolor{green!60!black}{\checkmark\ Done} \\
\hline
US-10 & Product CRUD (admin) & \textcolor{green!60!black}{\checkmark\ Done} \\
\hline
\end{tabular}
\end{table}

\subsection{Challenges}

\begin{itemize}
  \item \textbf{Checkout atomicity}: ensuring that stock decrement and order
        creation are inseparable required introducing an abstract
        \texttt{ITransactionManager}, injectable into the use case without
        coupling the domain layer to Drizzle.
  \item \textbf{JSONB filtering}: array columns (\texttt{skinType},
        \texttt{ingredients}) stored as \texttt{JSONB} require the PostgreSQL
        \texttt{@>} operator, not natively exposed by Drizzle; a typed raw SQL
        fragment was used.
  \item \textbf{Cart synchronization}: when logging in from a second device,
        the local (unauthenticated) cart must be merged with the server cart ---
        the merge logic was implemented in \texttt{CartUseCase}.
\end{itemize}

\section*{Conclusion}
This first iteration laid the foundations of the \textbf{LUMINA} e-commerce core.
All planned user stories were delivered, including the atomic checkout transaction
and hardened authentication (OTP + 2FA).
Iteration 2 will build on these foundations to add the artificial intelligence
layer: RAG chatbot and personalized recommendation engine.
